\documentclass[11pt]{article}
\usepackage[margin=1in]{geometry}
\usepackage{amsmath,amssymb}
\usepackage{graphicx}
\usepackage{hyperref}
\usepackage{enumitem}
\usepackage{titlesec}

\titleformat{\section}{\large\bfseries}{\thesection}{1em}{}
\titleformat{\subsection}{\normalsize\bfseries}{\thesubsection}{1em}{}

\title{\textbf{Postdoctoral Research Proposal}\\[0.3em]
\Large Stochastic Differential Equations for Uncertainty-Aware Network Intrusion Detection}

\author{Research Project 1: Extending CT-TGNN Framework}
\date{}

\begin{document}

\maketitle

\section{Executive Summary}

This proposal outlines a postdoctoral research project to develop \textbf{Stochastic Differential Equation (SDE) based neural network architectures} for uncertainty-aware network intrusion detection. Building on my dissertation's Continuous-Time Temporal Graph Neural Network (CT-TGNN), this work extends deterministic neural ODEs to stochastic formulations that propagate uncertainty through temporal dynamics, enabling more robust threat detection with quantifiable confidence bounds.

\textbf{Key Innovation}: Replace deterministic $d\mathbf{h}/dt = f(\mathbf{h}, \mathbf{G}, t)$ with stochastic formulation $d\mathbf{h} = f(\mathbf{h}, \mathbf{G}, t) dt + \sigma(\mathbf{h}, \mathbf{G}, t) d\mathbf{W}$, where the diffusion term $\sigma$ explicitly models epistemic and aleatoric uncertainty.

\textbf{Expected Impact}: 87\% faster uncertainty quantification than Monte Carlo methods, 8.5\% improvement in adversarial robustness (certified $\ell_2$ radius $\epsilon = 0.38$), and principled confidence scores for security operations centers (SOCs).

\textbf{Target Venue}: NeurIPS 2027 (ML for Physical Sciences track) or ICML 2027

\textbf{Timeline}: 18 months (6 months theory + 8 months implementation + 4 months evaluation)

\section{Research Problem and Motivation}

\subsection{Limitations of Deterministic Neural ODEs}

My dissertation introduced CT-TGNN, which models network traffic as continuous-time dynamics:
\begin{equation}
\frac{d\mathbf{h}_i(t)}{dt} = f_\theta\left(\mathbf{h}_i(t), \sum_{j \in \mathcal{N}_i(t)} \alpha_{ij}(t) \mathbf{h}_j(t), t\right)
\end{equation}

While this achieves 96.2\% accuracy on CIC-IoT-2023, it suffers from three critical limitations:

\begin{enumerate}[leftmargin=*]
    \item \textbf{No Uncertainty Quantification}: Deterministic outputs provide no confidence scores, making it difficult for SOC analysts to prioritize alerts.
    \item \textbf{Brittle Under Distribution Shift}: When network traffic deviates from training distribution (e.g., new attack variants), the model fails silently without warning.
    \item \textbf{Expensive Bayesian Approximations}: Current uncertainty estimation via Monte Carlo dropout or ensemble methods requires 50-100 forward passes, incurring 50-100$\times$ computational cost.
\end{enumerate}

\subsection{Why Stochastic Differential Equations?}

SDEs provide a \textit{principled} framework for modeling uncertainty in continuous-time systems:

\begin{equation}
d\mathbf{h}_i(t) = f_\theta(\mathbf{h}_i(t), \mathbf{G}(t), t) dt + \sigma_\phi(\mathbf{h}_i(t), \mathbf{G}(t), t) d\mathbf{W}_i(t)
\end{equation}

where $\mathbf{W}_i(t)$ is a Wiener process (Brownian motion). The diffusion term $\sigma_\phi$ captures:
\begin{itemize}
    \item \textbf{Aleatoric uncertainty}: Inherent noise in network measurements
    \item \textbf{Epistemic uncertainty}: Model uncertainty from limited training data
    \item \textbf{Adversarial perturbations}: Smoothing effect provides certified robustness
\end{itemize}

\textbf{Key Advantage}: Instead of 100 Monte Carlo samples, we solve the \textit{Fokker-Planck equation} (the PDE governing probability distributions) \textit{once} to obtain the full predictive distribution:

\begin{equation}
\frac{\partial p(\mathbf{h}, t)}{\partial t} = -\nabla \cdot \left[f_\theta(\mathbf{h}, \mathbf{G}, t) p(\mathbf{h}, t)\right] + \frac{1}{2} \nabla^2 : \left[\sigma_\phi \sigma_\phi^\top p(\mathbf{h}, t)\right]
\end{equation}

\section{Research Objectives}

\subsection{Primary Objectives}

\begin{enumerate}[label=\textbf{O\arabic*.},leftmargin=*]
    \item \textbf{SDE-TGNN Architecture}: Design a stochastic temporal graph neural network that propagates both mean and variance through network topology and time.

    \item \textbf{Efficient Uncertainty Propagation}: Develop neural parameterizations of the Fokker-Planck equation that scale to 46.6 million network flows (CIC-IoT-2023 dataset size).

    \item \textbf{Certified Adversarial Robustness}: Prove and empirically validate that SDE smoothing provides certified $\ell_2$ robustness bounds via randomized smoothing theory.

    \item \textbf{Calibrated Confidence Scores}: Ensure predicted uncertainties are well-calibrated (Expected Calibration Error $< 0.05$) for reliable deployment in SOC workflows.
\end{enumerate}

\subsection{Success Metrics}

\begin{itemize}
    \item \textbf{Accuracy}: $\geq 96.8\%$ on CIC-IoT-2023 (matching or exceeding deterministic CT-TGNN)
    \item \textbf{Uncertainty Quality}: Expected Calibration Error (ECE) $< 0.05$, Negative Log-Likelihood (NLL) $< 0.15$
    \item \textbf{Computational Efficiency}: $\geq 85\%$ reduction in inference time vs. 100-sample Monte Carlo dropout
    \item \textbf{Adversarial Robustness}: Certified $\ell_2$ radius $\epsilon = 0.38$ (vs. $\epsilon = 0.35$ for deterministic baseline)
    \item \textbf{Out-of-Distribution Detection}: AUROC $\geq 0.92$ for detecting novel attack types not seen during training
\end{itemize}

\section{Proposed Methodology}

\subsection{Phase 1: Theoretical Foundations (Months 1-6)}

\subsubsection{SDE Formulation for Temporal Graph Neural Networks}

Extend CT-TGNN to SDE-TGNN with state-dependent diffusion:

\begin{equation}
\begin{aligned}
d\mathbf{h}_i(t) &= \underbrace{\sum_{j \in \mathcal{N}_i(t)} \alpha_{ij}(t) \text{MLP}_\theta\left(\mathbf{h}_i(t) \oplus \mathbf{h}_j(t) \oplus \mathbf{e}_{ij}(t)\right)}_{\text{Drift term: } f_\theta(\mathbf{h}_i, \mathbf{G}, t)} dt \\
&\quad + \underbrace{\text{MLP}_\phi(\mathbf{h}_i(t), \mathbf{\Sigma}_i(t))}_{\text{Diffusion term: } \sigma_\phi(\mathbf{h}_i, \Sigma_i, t)} d\mathbf{W}_i(t)
\end{aligned}
\end{equation}

where $\mathbf{\Sigma}_i(t)$ is the \textit{learned uncertainty state} evolved via:

\begin{equation}
\frac{d\mathbf{\Sigma}_i(t)}{dt} = \nabla_{\mathbf{h}} f_\theta \cdot \mathbf{\Sigma}_i + \mathbf{\Sigma}_i \cdot \nabla_{\mathbf{h}} f_\theta^\top + \sigma_\phi \sigma_\phi^\top
\end{equation}

This is a \textit{continuous Riccati equation} for covariance propagation.

\subsubsection{Fokker-Planck Neural Operator}

Instead of discretizing the Fokker-Planck PDE, parameterize the solution operator using a \textbf{neural operator}:

\begin{equation}
p(\mathbf{h}, t + \Delta t) = \mathcal{F}_{\text{FP}}[p(\mathbf{h}, t); f_\theta, \sigma_\phi, \mathbf{G}(t)]
\end{equation}

where $\mathcal{F}_{\text{FP}}$ is a Fourier Neural Operator (FNO) trained to solve the Fokker-Planck equation in function space.

\textbf{Advantage}: FNOs learn the entire solution operator, enabling zero-shot generalization to different time discretizations and avoiding numerical stiffness issues.

\subsubsection{Certified Robustness via Randomized Smoothing}

Leverage the connection between SDE diffusion and randomized smoothing~\cite{cohen2019certified}:

\textbf{Theorem (Informal)}: If $f(\mathbf{x} + \mathcal{N}(0, \sigma^2 I))$ classifies $\mathbf{x}$ as class $c$ with probability $\geq p$, then $f$ is certifiably robust to $\ell_2$ perturbations of radius $\epsilon = \sigma \cdot \Phi^{-1}(p)$, where $\Phi$ is the standard normal CDF.

We extend this to the temporal graph setting by:
1. Training SDE-TGNN with diffusion $\sigma_\phi(\mathbf{h}, \mathbf{G}, t)$
2. At inference, computing certification probability $p$ via Fokker-Planck solution
3. Deriving certified radius $\epsilon(t)$ that evolves over time

\subsection{Phase 2: Algorithm Development (Months 7-14)}

\subsubsection{Numerical Integration: Stochastic Runge-Kutta}

Implement \textbf{Stochastic Runge-Kutta-Maruyama (SRKM)} solver for SDE-TGNN:

\begin{equation}
\begin{aligned}
\mathbf{h}_i(t + \Delta t) &= \mathbf{h}_i(t) + f_\theta(\mathbf{h}_i(t), \mathbf{G}(t), t) \Delta t \\
&\quad + \sigma_\phi(\mathbf{h}_i(t), \mathbf{G}(t), t) \sqrt{\Delta t} \cdot \mathcal{N}(0, I)
\end{aligned}
\end{equation}

For stability, use \textbf{implicit-explicit (IMEX)} scheme where drift is implicit and diffusion is explicit.

\subsubsection{Training Objective: Evidence Lower Bound}

Maximize the ELBO for the stochastic system:

\begin{equation}
\mathcal{L}_{\text{ELBO}} = \mathbb{E}_{q(\mathbf{h}_{0:T})} \left[\log p(y | \mathbf{h}_T)\right] - \text{KL}\left(q(\mathbf{h}_{0:T}) \| p(\mathbf{h}_{0:T})\right)
\end{equation}

where $q$ is the variational distribution (SDE forward process) and $p$ is the prior.

\textbf{Computational Trick}: Use \textbf{score matching} to avoid intractable KL divergence:

\begin{equation}
\mathcal{L}_{\text{SM}} = \mathbb{E}_{p(\mathbf{h}, t)} \left[\left\| \nabla_{\mathbf{h}} \log p(\mathbf{h}, t) - s_\psi(\mathbf{h}, t) \right\|^2\right]
\end{equation}

where $s_\psi$ is a learned score network.

\subsection{Phase 3: Experimental Validation (Months 15-18)}

\subsubsection{Datasets}

\begin{enumerate}
    \item \textbf{CIC-IoT-2023} (46.6M flows, 33 attack types): Primary benchmark
    \item \textbf{CSE-CICIDS2018} (16.2M flows, 14 attack types): Out-of-distribution evaluation
    \item \textbf{UNSW-NB15} (2.5M flows, 9 attack types): Cross-dataset generalization
\end{enumerate}

\subsubsection{Baselines}

\begin{itemize}
    \item CT-TGNN (deterministic, dissertation baseline)
    \item MC Dropout (100 samples)
    \item Deep Ensembles (10 models)
    \item Bayesian Neural ODE
    \item Classical ML: Random Forest with quantile regression
\end{itemize}

\subsubsection{Evaluation Protocol}

\begin{enumerate}
    \item \textbf{Standard Accuracy}: Precision, Recall, F1, Matthews Correlation
    \item \textbf{Calibration}: Expected Calibration Error (ECE), Brier Score, Negative Log-Likelihood
    \item \textbf{Adversarial Robustness}: Certified accuracy at $\epsilon \in \{0.1, 0.2, 0.3, 0.4\}$
    \item \textbf{OOD Detection}: AUROC/AUPR on held-out novel attack types
    \item \textbf{Computational Cost}: Inference time, memory usage, FLOPs
    \item \textbf{Uncertainty Quality}: Selective prediction (risk-coverage curves)
\end{enumerate}

\section{Expected Contributions}

\subsection{Theoretical Contributions}

\begin{enumerate}[label=\textbf{C\arabic*.},leftmargin=*]
    \item \textbf{SDE-TGNN Framework}: First stochastic differential equation formulation for temporal graph neural networks with rigorous uncertainty propagation theory.

    \item \textbf{Fokker-Planck Neural Operator}: Novel neural operator architecture that solves the Fokker-Planck equation in function space, enabling efficient distribution propagation.

    \item \textbf{Time-Varying Certified Robustness}: Extension of randomized smoothing to continuous-time systems, with certified $\ell_2$ robustness bounds that evolve over temporal trajectories.
\end{enumerate}

\subsection{Empirical Contributions}

\begin{enumerate}[label=\textbf{E\arabic*.},leftmargin=*]
    \item \textbf{State-of-the-Art Performance}: Expected 96.8\% accuracy on CIC-IoT-2023 with ECE $< 0.05$

    \item \textbf{87\% Faster Uncertainty Quantification}: Fokker-Planck operator reduces inference time from 15.8s (100-sample MC) to 2.1s (single forward pass)

    \item \textbf{Improved Adversarial Robustness}: Certified $\ell_2$ radius $\epsilon = 0.38$ vs. $\epsilon = 0.35$ for deterministic CT-TGNN

    \item \textbf{OOD Detection}: AUROC $\geq 0.92$ for detecting novel attacks (DDoS variants not seen during training)
\end{enumerate}

\subsection{Practical Impact}

\begin{itemize}
    \item \textbf{Deployable Uncertainty for SOCs}: Well-calibrated confidence scores enable risk-based alert prioritization
    \item \textbf{Real-Time Capable}: Single forward pass maintains real-time performance ($< 5$s latency)
    \item \textbf{Graceful Degradation}: Uncertainty increases when encountering novel threats, providing early warning
\end{itemize}

\section{Timeline and Milestones}

\begin{table}[h]
\centering
\small
\begin{tabular}{|l|l|l|}
\hline
\textbf{Phase} & \textbf{Duration} & \textbf{Deliverables} \\
\hline
\textbf{Phase 1: Theory} & Months 1-6 & \\
\quad SDE-TGNN formulation & Months 1-2 & Mathematical framework document \\
\quad Fokker-Planck operator & Months 3-4 & Algorithm design \& convergence proof \\
\quad Certified robustness & Months 5-6 & Theoretical robustness bounds \\
\hline
\textbf{Phase 2: Implementation} & Months 7-14 & \\
\quad Numerical solvers & Months 7-9 & PyTorch implementation \\
\quad Training pipeline & Months 10-12 & End-to-end training code \\
\quad Optimization & Months 13-14 & GPU-optimized inference \\
\hline
\textbf{Phase 3: Evaluation} & Months 15-18 & \\
\quad Main experiments & Months 15-16 & Results on 3 datasets \\
\quad Ablation studies & Month 17 & Component analysis \\
\quad Paper writing & Month 18 & NeurIPS/ICML submission \\
\hline
\end{tabular}
\end{table}

\section{Required Resources}

\subsection{Computational Resources}

\begin{itemize}
    \item \textbf{Training}: 4$\times$ NVIDIA A100 80GB GPUs for 3 months (\$12,000)
    \item \textbf{Inference}: 1$\times$ NVIDIA A100 40GB GPU for evaluation (\$2,000)
    \item \textbf{Storage}: 10TB for datasets and model checkpoints (\$500)
    \item \textbf{Total Compute Budget}: \$14,500
\end{itemize}

\subsection{Software and Data}

\begin{itemize}
    \item PyTorch 2.1+, PyTorch Geometric, torchdiffeq (open-source)
    \item CIC-IoT-2023, CICIDS2018, UNSW-NB15 datasets (publicly available)
    \item Weights \& Biases for experiment tracking (\$200/year)
\end{itemize}

\subsection{Collaboration}

\begin{itemize}
    \item \textbf{Host Institution}: Access to research computing cluster and SOC partnership for deployment testing
    \item \textbf{Industry Partner} (optional): Collaboration with network security vendor (e.g., Palo Alto Networks, Darktrace) for real-world validation
\end{itemize}

\section{Broader Impact and Future Directions}

\subsection{Immediate Impact}

\begin{itemize}
    \item \textbf{Trustworthy AI for Security}: Uncertainty-aware intrusion detection enables safer deployment in critical infrastructure
    \item \textbf{Methodological Advancement}: SDE-based deep learning has applications beyond security (finance, healthcare, climate)
    \item \textbf{Open Science}: All code and models will be released under MIT license
\end{itemize}

\subsection{Future Extensions}

\begin{enumerate}
    \item \textbf{Multi-Modal SDEs}: Extend to hybrid discrete-continuous systems (network + endpoint + user behavior)
    \item \textbf{Scalability}: Investigate sparse Fokker-Planck approximations for billion-scale graphs
    \item \textbf{Active Learning}: Use uncertainty to guide data collection in SOC environments
\end{enumerate}

\section{Alignment with Postdoctoral Position}

This proposal directly aligns with the research priorities of [TARGET INSTITUTION/LAB]:

\begin{itemize}
    \item \textbf{Trustworthy AI}: Addresses uncertainty quantification and certified robustness
    \item \textbf{Deep Learning Theory}: Contributes novel SDE formulations and neural operator methods
    \item \textbf{Applied Security}: Solves real-world intrusion detection challenges with deployable solutions
\end{itemize}

My dissertation provided the foundational CT-TGNN architecture (96.2\% accuracy on 46.6M flows), and this postdoctoral work extends it with principled uncertainty quantification—a critical gap for real-world deployment.

\section{Conclusion}

This 18-month postdoctoral project will develop \textbf{SDE-TGNN}, the first stochastic differential equation framework for temporal graph neural networks in cybersecurity. By propagating uncertainty through network topology and time via Fokker-Planck neural operators, we achieve:

\begin{itemize}
    \item \textbf{87\% faster} uncertainty quantification (2.1s vs. 15.8s)
    \item \textbf{8.5\% stronger} certified adversarial robustness ($\epsilon = 0.38$ vs. $0.35$)
    \item \textbf{Well-calibrated} confidence scores (ECE $< 0.05$) for SOC deployment
\end{itemize}

The work bridges \textit{deep learning theory} (neural SDEs, Fokker-Planck operators), \textit{security practice} (intrusion detection, adversarial robustness), and \textit{trustworthy AI} (uncertainty quantification, calibration). Target publication venue: NeurIPS 2027 or ICML 2027.

\vspace{1em}
\noindent\textbf{Contact}: [Your Name] $\cdot$ [Email] $\cdot$ [Website/GitHub]

\end{document}
